\begin{lemma}[Equivalent criterion for stability]
Let $\varphi$ and $\pi$ be proper partial $k$-edge-colorings which leave
exactly the same distinguished edge $e_0$ uncolored.  Then $\pi$ is
$(T,C,\varphi)$-stable if and only if, for every edge identity $f$ incident
with $T$,
\[
 \varphi(f)\in\overline\varphi(T)\cup C
 \quad\hbox{or}\quad
 \pi(f)\in\overline\varphi(T)\cup C
 \quad\Longrightarrow\quad
 \pi(f)=\varphi(f).
\]
Here a displayed membership requires the coloring value to be a color, not
the uncolored value.
\end{lemma}
\begin{proof}
Write $(i')$ for the displayed condition and use the definitions of
tree-incidence, tree-missing colors, and stability in
\noderef{TreeIncidentEdge}, \noderef{TreeMissingColors}, and
\noderef{StableColoring}.

Assume first $(i')$.  Its case in which
$\varphi(f)\in\overline\varphi(T)\cup C$ is exactly condition (i) of
stability.  We prove condition (ii).  Fix $v\in V(T)$.  Because both
colorings leave exactly the same edge $e_0$ uncolored,
\noderef{IsOneEdgeUncolored} says that each coloring colors every edge
identity incident with $v$ except $e_0$ when $e_0$ is incident with $v$.
By \noderef{Degree}, $d_G(v)$ is the number of incident edge identities.
Thus each coloring colors $d_G(v)$ incident edges if $e_0$ is not incident
with $v$, and $d_G(v)-1$ incident edges if it is.  By properness,
\noderef{IsProperEdgeColoring} says that distinct incident colored edges
have distinct colors.  Therefore the definition in \noderef{PresentColors}
identifies each present color at $v$ with exactly one of those colored
incident edges.  The two present-color sets at $v$ consequently have equal
cardinality, and so do their complements in $[k]$, the missing-color sets
of \noderef{MissingColors}.

If those two missing-color sets were unequal, equal cardinality would give
a color $c$ which is missing under $\varphi$ but present under $\pi$ at
$v$.  By \noderef{PresentColors}, the fact that $c$ is present under
$\pi$ supplies an edge identity $f$ incident with $v$ such that
$\pi(f)=c$.  The edge $f$ is incident with $T$, and
$c\in\overline\varphi(v)$ implies
$c\in\overline\varphi(T)$.  Condition $(i')$ therefore gives
$\varphi(f)=\pi(f)=c$, contradicting that $c$ is missing under $\varphi$
at $v$.  Thus the missing-color sets agree at every vertex of $T$, proving
condition (ii) and hence stability.

Conversely, assume that $\pi$ is stable.  Condition (i) already treats the
case in which $\varphi(f)$ lies in the specified set.  Let $f$ be incident
with $T$, suppose instead that $\pi(f)=c$ for a color
$c\in\overline\varphi(T)\cup C$, and suppose for contradiction that
$\varphi(f)\ne c$.  Choose an end $v\in V(T)$ of $f$.  Stability condition
(ii) makes the present-color sets of $\varphi$ and $\pi$ equal at $v$.
Since $c$ is present under $\pi$, it is present under $\varphi$, so some
edge $g$ incident with $v$ satisfies $\varphi(g)=c$, by
\noderef{PresentColors}.  Our supposition makes $g\ne f$.  The edge $g$ is
incident with $T$ and its $\varphi$-color $c$ lies in
$\overline\varphi(T)\cup C$; stability condition (i) therefore gives
$\pi(g)=\varphi(g)=c$.  The distinct edges $f$ and $g$ share the vertex
$v$, so they are adjacent by \noderef{AdjacentEdges}.  They now both have
color $c$ under $\pi$, contradicting the properness condition of
\noderef{IsProperEdgeColoring}.  Hence $\pi(f)=\varphi(f)$.  Combining the
two cases proves $(i')$ and the equivalence.
\end{proof}
